\section{Conclusiones}

En esta práctica se implementó el algoritmo de Gale--Shapley para resolver el Problema de Emparejamiento Estable y se evaluó experimentalmente su comportamiento temporal para distintos tamaños de entrada. Con ello, se cumplió el objetivo principal de analizar cómo evoluciona el costo computacional del algoritmo conforme aumenta el valor de $N$.

Los resultados obtenidos muestran que el tiempo de ejecución crece de manera aproximadamente cuadrática con respecto a $N$. El ajuste realizado en escala log-log arrojó una pendiente cercana a $1.96$, valor consistente con la complejidad teórica $O(N^2)$ del algoritmo en el peor caso. Por lo tanto, el análisis empírico concuerda con lo establecido en la literatura para el algoritmo de Gale--Shapley.

Las variaciones observadas entre los datos experimentales y el modelo ideal pueden explicarse por factores externos al algoritmo, como el entorno de ejecución, la administración de recursos del sistema y la variabilidad propia de cada instancia generada. No obstante, dichas diferencias son menores frente a la tendencia general identificada en la tabla y en las gráficas.

En conclusión, la práctica no solo permitió implementar correctamente el procedimiento de emparejamiento estable, sino también comprobar mediante evidencia experimental que su crecimiento temporal es coherente con un comportamiento cuadrático. De esta manera, los objetivos planteados se consideran alcanzados y los resultados obtenidos guardan concordancia con el análisis presentado.
